\begin{abstract}
\noindent\textbf{Background.} Machine-learning studies often report near-perfect accuracy for detecting major
depressive disorder (MDD) from EEG. Many of them, however, split trials rather than participants between training and
test sets, and few compare against non-EEG baselines.

\noindent\textbf{Methods.} We re-analysed the 128-channel dot-probe ERP recordings of the MODMA dataset (24 MDD, 29 HC)
from the raw files with a pre-specified, version-controlled plan. Preprocessing used PREP bad-channel detection and
ICA with ICLabel, and participants were excluded by fixed quality criteria. We evaluated 18 model configurations:
\begin{itemize}
  \item a demographic baseline;
  \item linear, kernel and tree models on 288 handcrafted features (spectral, frontal alpha asymmetry, ERP, wPLI
        connectivity, complexity);
  \item Riemannian models;
  \item three deep networks on full epochs.
\end{itemize}
All models used nested cross-validation over participants with identical folds. Results carry participant-level
bootstrap confidence intervals, and all but the deep networks carry permutation tests of the full pipeline.

\noindent\textbf{Results.} In \NIncluded{} participants, the pre-specified primary model (L2 logistic regression on all
features) reached AUC \PrimaryLrLAllAuc{} (95\% CI \PrimaryLrLAllAucLo--\PrimaryLrLAllAucHi, $p=\PrimaryLrLAllPperm$).
No EEG model exceeded a baseline using only age, sex and education (AUC \DemographicsLrAuc). The best secondary model
(AUC \LrLAllAuc) did not survive multiplicity correction, and nested model selection gave AUC
\NestedmodelselectionAllAuc. Evaluated with an invalid trial-wise split, the same tangent-space model achieved apparently perfect
participant-level classification (AUC \InflTsLrTrialSplitSubjectAuc), versus \InflTsLrSubjectSplitSubjectAuc{} when
participants were held out. The null
result held when all 53 participants were included and when the analysis was restricted to a period in which both
groups were recorded. Post-hoc exploratory analyses did not change it either:
\begin{itemize}
  \item EEG added nothing to demographics (AUC \AudDemoPlusEEGAllAuc);
  \item EEG with demographics regressed out inside each fold gave AUC \AudEEGAlldeconfoundedAuc;
  \item compact theory-driven features, a subject-level Riemannian model and reaction times were all at chance;
  \item the only borderline signal came from a random forest on raw ERP waveforms ($p=\AudRawRfP$, not significant after correction), which we propose as a hypothesis for independent testing.
\end{itemize}
Complexity and high-frequency features tracked a raw-data EMG index. With 18 patients, only large effects were
detectable.

\noindent\textbf{Conclusions.} This study contributes a leakage-free reference benchmark for MODMA and identifies
the mechanism behind inflated accuracies. A trial-level classifier recognised \emph{which} participant produced a
trial with \AudIdentAcc\% accuracy, so trial-wise splits reward person recognition rather than diagnosis. Evaluated
on unseen participants, task EEG in this dataset carried no detectable diagnostic signal beyond a demographic baseline
that reflects recruitment differences. Residual muscle activity is a concrete, checkable confound for complexity and
high-frequency features. Even an ideal score with a modest true effect (AUC 0.65) would need about
\AudNNeedSixtyFive{} participants to be detected reliably. We provide a reporting checklist and open code so that
future EEG studies of depression can be evaluated, and compared, on the same footing.
\end{abstract}
